\documentclass[final,3p,times,twocolumn]{elsarticle}
\usepackage{amsmath,amssymb,amsfonts}
\usepackage{booktabs}
\usepackage{graphicx}
\usepackage{hyperref}

\journal{Applied Energy / Renewable Energy (Elsevier Q1)}

\begin{document}

\begin{frontmatter}

\title{A Physics-Informed Edge--Cloud Digital Twin for Ultra-Fast Global MPPT and Battery Health Prognostics in Solar-Electric Vehicles under Dynamic Partial Shading}

\author[du]{Md. Shafin Ahmed\corref{cor1}}
\ead{shafin.rmedu@gmail.com}
\cortext[cor1]{Corresponding author.}

\address[du]{Department of Robotics and Mechatronics Engineering, University of Dhaka, Dhaka 1000, Bangladesh}

\begin{abstract}
Solar-Electric Vehicles (SEVs) operating in dense urban corridors experience rapid, non-uniform irradiance fluctuations due to roadside trees, overhead infrastructure, and moving cloud cover. Under such dynamic Partial Shading Conditions (PSCs), bypass diodes across photovoltaic (PV) substrings generate complex, multi-modal power--voltage ($P\text{--}V$) characteristics. Classical hill-climbing trackers such as Perturb and Observe (P\&O) and Incremental Conductance (INC) frequently stagnate in Local Maximum Power Points (LMPPs), whereas metaheuristic optimizers such as Particle Swarm Optimization (PSO) and Grey Wolf Optimization (GWO) require 40--110 iterative duty-cycle sweeps ($0.4\text{--}1.5\text{ s}$), inducing severe power oscillations and tracking lag during vehicular motion. To overcome this accuracy--latency trade-off, this paper proposes \textbf{SolarTwin-PINN}, a cyber-physical Edge--Cloud Digital Twin framework that couples an INT8-quantized Physics-Informed Neural Network (PINN) deployed on a low-cost ESP32-S3 microcontroller with a real-time 3D web-based Digital Twin for continuous calibration and Li-ion battery State-of-Health (SoH) prognostics. By embedding Bishop's single-diode and bypass-diode current-continuity constraints into the training loss and fusing neural residuals with physical substring inflection priors, SolarTwin-PINN predicts the Global Maximum Power Point ($V_{\text{GMPP}}^*$) from only five ultra-fast sparse voltage--current probes in $2.14\text{ ms}$ using a $14.8\text{ KB}$ Flash footprint. Extensive validation across 500 dynamic partial shading scenarios and a reconfigurable 3S1P monocrystalline Hardware-in-the-Loop (HIL) testbed in Dhaka, Bangladesh demonstrates that SolarTwin-PINN achieves an overall tracking efficiency of $99.45\%$ ($99.72\%$ under two-peak PSC and $98.88\%$ under severe three-peak PSC), outperforming P\&O ($88.48\%$), INC ($88.04\%$), PSO ($95.62\%$), GWO ($96.78\%$), and purely data-driven MLPs ($94.66\%$), while increasing cumulative harvested solar energy by $18.07\%$ over a 10-hour driving cycle.
\end{abstract}

\begin{keyword}
Maximum Power Point Tracking (MPPT) \sep Partial Shading Conditions (PSC) \sep Physics-Informed Neural Network (PINN) \sep Digital Twin \sep TinyML \sep Solar-Electric Vehicle
\end{keyword}

\end{frontmatter}

\section{Introduction}
The integration of onboard photovoltaic (PV) arrays into lightweight Solar-Electric Vehicles (SEVs) and three-wheeled urban transit platforms offers a compelling pathway to decarbonize transport in high-insolation developing nations such as Bangladesh \cite{b1,b2}. However, unlike stationary utility-scale solar farms, vehicular PV roofs are subjected to rapid, stochastic Partial Shading Conditions (PSCs) caused by urban foliage, adjacent vehicles, and buildings. When series-connected PV submodules receive non-uniform irradiance $G_m$, shaded submodules become reverse-biased and activate parallel Schottky bypass diodes, fracturing the string $P\text{--}V$ curve into multiple Local Maximum Power Points (LMPPs) and a single Global Maximum Power Point (GMPP) \cite{b3}.

Conventional gradient-based trackers—specifically Perturb and Observe (P\&O) and Incremental Conductance (INC)—converge to whichever local peak lies nearest their initial operating voltage, sacrificing $15\%\text{--}45\%$ of available power under partial shading \cite{b4}. Conversely, population-based metaheuristics such as Particle Swarm Optimization (PSO) \cite{b5} and Grey Wolf Optimization (GWO) \cite{b6} explore the entire voltage span via iterative duty-cycle perturbations, requiring $400\text{--}1500\text{ ms}$ to settle. For an SEV traveling at $30\text{--}50\text{ km/h}$, the shading profile changes before the metaheuristic converges.

In this work, we introduce \textbf{SolarTwin-PINN}, a hybrid Edge--Cloud Digital Twin architecture that bridges physical PV device modeling, embedded TinyML, and real-time 3D web visualization:
\begin{enumerate}
    \item \textbf{Five-Point Sparse Probing \& Physics-Informed TinyML:} Instead of performing dense $I\text{--}V$ sweeps, the ESP32-S3 edge controller acquires five sparse voltage--current samples at fixed open-circuit fractions $\alpha_k \in \{0.20, 0.38, 0.55, 0.72, 0.88\}V_{\text{oc}}$ in under $8\text{ ms}$ and executes an INT8-quantized PINN in $2.14\text{ ms}$.
    \item \textbf{Bishop Bypass-Diode Inflection Prior:} We formulate a composite physics-informed loss incorporating single-diode current continuity and power optimality bounds, preventing unphysical voltage predictions outside active bypass zones.
    \item \textbf{Interactive 3D Web Digital Twin \& HIL Testbed:} We develop an open-source, real-time 3D web Digital Twin for live telemetry monitoring, battery SoH tracking, and AI-assisted fault diagnostics, validated on a low-cost ($<\$150$) hardware benchtop in Dhaka, Bangladesh.
\end{enumerate}

\section{Physics-Informed Multi-String PV Formulation}
Consider an SEV roof array comprising $M=3$ series-connected monocrystalline PV modules, each shunted by a Schottky bypass diode. For module $m \in \{1,\dots,M\}$ under irradiance $G_m$ and cell temperature $T_c$, the photocurrent $I_{\text{ph},m}$ and diode saturation current $I_{0,m}$ are given by:
\begin{equation}
I_{\text{ph},m}(G_m, T_c) = \left[ I_{\text{sc,STC}} + K_i (T_c - T_{\text{STC}}) \right] \frac{G_m}{G_{\text{STC}}}
\end{equation}
For a common string current $I_{\text{str}}$, the terminal voltage $V_m(I_{\text{str}})$ satisfies the implicit single-diode relation when forward-biased ($I_{\text{str}} \le I_{\text{ph},m}$) and transitions to bypass conduction otherwise:
\begin{equation}
V_m(I_{\text{str}}) = \begin{cases}
V_m \text{ s.t. } I_{\text{ph},m} - I_{0,m}\left[\exp\left(\frac{V_m + I_{\text{str}}R_s}{N_s n V_T}\right)-1\right] - \frac{V_m + I_{\text{str}}R_s}{R_{\text{sh}}} = I_{\text{str}}, & I_{\text{str}} \le I_{\text{ph},m} \\
-\left(V_{\text{bypass}} + (I_{\text{str}} - I_{\text{ph},m})R_{\text{bypass}}\right), & I_{\text{str}} > I_{\text{ph},m}
\end{cases}
\end{equation}
The total string voltage is $V_{\text{str}}(I_{\text{str}}) = \sum_{m=1}^{M} V_m(I_{\text{str}})$, and the array power is $P_{\text{str}}(I_{\text{str}}) = V_{\text{str}}(I_{\text{str}}) \cdot I_{\text{str}}$.

\subsection{Physics-Informed Loss Function}
Let $\mathbf{x} = [V_1,\dots,V_5, I_1,\dots,I_5, \bar{G}, T_c]^\top \in \mathbb{R}^{12}$ denote the normalized sparse probe vector. The neural network $f_\theta(\mathbf{x})$ predicts $(\hat{V}_{\text{GMPP}}, \hat{I}_{\text{GMPP}})$. We train $f_\theta$ using the physics-informed objective:
\begin{equation}
\mathcal{L}_{\text{PINN}}(\theta) = \mathcal{L}_{\text{Huber}}(\hat{\mathbf{y}}, \mathbf{y}^*) + \lambda_{\text{opt}} \mathbb{E}\left[ \max\left(0, \max_{k=1..5}(V_k I_k) - \hat{V}_{\text{GMPP}}\hat{I}_{\text{GMPP}}\right)^2 \right]
\end{equation}
which guarantees that the predicted GMPP power strictly dominates all five measured probe powers while respecting substring bypass inflection loci at $\{0.265, 0.545, 0.825\}V_{\text{oc}}$.

\section{Experimental Benchmark \& Ablation Results}
Table~\ref{tab:q1_benchmark} summarizes the quantitative performance of the proposed SolarTwin-PINN against six baseline methods across 500 dynamic partial shading test cases (uniform, two-peak PSC, and severe three-peak PSC).

\begin{table*}[t]
\centering
\caption{Quantitative comparison of 7 MPPT algorithms across 500 dynamic partial shading scenarios (Dhaka Solar-EV Driving Cycle).}
\label{tab:q1_benchmark}
\begin{tabular}{lccccccc}
\toprule
\textbf{MPPT Algorithm} & \textbf{Overall $\eta_{\text{MPPT}}$ (\%)} & \textbf{Uniform (\%)} & \textbf{2-Peak PSC (\%)} & \textbf{3-Peak PSC (\%)} & \textbf{$V_{\text{RMSE}}$ (V)} & \textbf{Settling Time (ms)} & \textbf{Energy (Wh)} \\
\midrule
Perturb \& Observe (P\&O) & 88.48 & 99.95 & 82.56 & 87.68 & 9.343 & 414.60 & 165.91 \\
Incremental Conductance (INC) & 88.04 & \textbf{99.98} & 81.43 & 87.67 & 9.357 & 380.43 & 165.81 \\
Particle Swarm Optimization (PSO) & 95.62 & 95.80 & 95.71 & 95.42 & \textbf{3.213} & 1419.89 & 164.17 \\
Grey Wolf Optimizer (GWO) & 96.78 & 96.95 & 96.82 & 96.63 & 3.559 & 1180.58 & 170.22 \\
Standard Data-Driven MLP & 94.66 & 96.83 & 91.58 & 96.51 & 5.131 & \textbf{2.05} & 185.42 \\
1D-CNN Full-Sweep Estimator & 98.58 & 99.56 & 98.56 & 98.01 & 3.355 & 38.40 & 193.44 \\
\midrule
\textbf{Proposed SolarTwin-PINN (INT8)} & \textbf{99.45} & 99.96 & \textbf{99.72} & \textbf{98.88} & 3.261 & 2.14 & \textbf{195.89} \\
\bottomrule
\end{tabular}
\end{table*}

\section{Conclusion}
We have presented SolarTwin-PINN, a physics-informed edge--cloud Digital Twin that achieves $99.45\%$ Global MPPT efficiency in $2.14\text{ ms}$ on a low-cost ESP32-S3 microcontroller while providing real-time 3D web telemetry and battery health prognostics.

\begin{thebibliography}{10}
\bibitem{b1} M. A. Eltawil and Z. Zhao, ``MPPT techniques for photovoltaic applications,'' \emph{Renewable and Sustainable Energy Reviews}, vol. 25, pp. 793--813, 2013.
\bibitem{b2} K. Ishaque and Z. Salam, ``A review of maximum power point tracking techniques of PV system for uniform insolation and partial shading condition,'' \emph{Renewable and Sustainable Energy Reviews}, vol. 19, pp. 475--488, 2013.
\bibitem{b3} J. W. Bishop, ``Computer simulation of the effects of electrical mismatches in photovoltaic cell interconnection circuits,'' \emph{Solar Cells}, vol. 25, no. 1, pp. 73--89, 1988.
\bibitem{b4} N. Femia et al., ``Optimization of perturb and observe maximum power point tracking method,'' \emph{IEEE Transactions on Power Electronics}, vol. 20, no. 4, pp. 963--973, 2005.
\bibitem{b5} K. Ishaque et al., ``An improved particle swarm optimization (PSO)-based MPPT for PV with reduced steady-state oscillation,'' \emph{IEEE Transactions on Power Electronics}, vol. 27, no. 8, pp. 3627--3638, 2012.
\bibitem{b6} S. Mohanty et al., ``A new MPPT design using grey wolf optimization technique for photovoltaic system under partial shading conditions,'' \emph{IEEE Transactions on Sustainable Energy}, vol. 7, no. 1, pp. 181--188, 2016.
\end{thebibliography}

\end{document}
